The original data files contained both the data and information in text regarding the data. To clean
the files into a readable csv file without the text, we used a .sh cleaning script written by Florido
Paganelli. This script develops different versions of the data file but the one that we preferred
was a script with each line having the format of "year-month-day hour:minute:second temperature
quality". The relevant parts of the cleaner script for this data format is shown.
To start, using the script above, all the beginning lines with just text are completely. This works
by reading the first lines of actual data with "STARTLINE" and directs the data to an output data
file with ">".

echo "Removing the first $STARTLINE lines, result in clean1_${DATAFILE}"
tail -n +$STARTLINE original_${DATAFILE} > clean1_${DATAFILE}

Then it cuts the text that is written next to data with the code below by using "cut" to select
only the first 5 columns and redirects that to an output data file:

echo "Selecting only relevant columns, result in clean2_${DATAFILE}"
cut -d';' -f 1,2,3,4,5 clean1_${DATAFILE} > clean2_${DATAFILE}

Lastly, it removed the semicolons that separate the data with "sed" and the semicolons were
replaced with spaces and created the final cleaned csv file. The substitution is given with "’s/;/
/g’":

echo "Substituting the ; with spaces, result in baredata_${DATAFILE}"
sed 's/;/ /g' clean2_${DATAFILE} > baredata_${DATAFILE}

The final version of the cleaned data file is of the desired format:
1780-01-01 07:00:00 -0.5 Y
1780-01-01 13:00:00 1.0 Y
1780-01-01 21:00:00 0.8 Y
1780-01-02 21:00:00 -3.0 Y
1780-01-03 13:00:00 1.3 Y
1780-01-03 21:00:00 -6.8 Y
1780-01-04 13:00:00 -9.0 Y
1780-01-05 13:00:00 -7.0 Y
1780-01-05 21:00:00 -10.0 Y
1780-01-06 21:00:00 -14.0 Y

Once we had the cleaned data file, the plan was to create a ROOT file with a branch for the
variables: year, month, day, hour, minute, second, temp, and quality. This was done through a
Measurement.cxx file and a DataExtraction.cxx file. The outlines for each script was used from
example scripts by Oxana Smirnova and Einar Elen.
The first script, Measurement.cxx was used to define how a measurement is accessed and stored:

//Measurement.cxx
#include "Measurement.h"
//Constructor to initialize variables
Measurement::Measurement(int year, int month, int day, int hour, int minute, int second, double temp
: year{year}, month{month}, day{day}, hour{hour}, minute{minute}, second{second}, temp{temp}, qu
}
//getter functions to get the value of an object and be stored in the member value (the returned var
int Measurement::get_year() const{
return year;
}
int Measurement::get_month() const{
return month;
}
int Measurement::get_day() const{
return day;
}
int Measurement::get_hour() const{
return hour;
}
int Measurement::get_minute() const{
return minute;
}
int Measurement::get_second() const{
return second;
}
double Measurement::get_temp() const{
return temp;
}
char Measurement::get_quality() const{
return quality;
}
bool Measurement::is_valid() const{
return true;
}

The DataExtraction.cxx code is able to then read the csv file and use Measurement.cxx for the
objects defined:

//DataExtraction.cxx
#include "DataExtraction.h"
//Measurment is a class defined in Measurement.h
//define the function read_measurements with the name given by the filename parameter
std::vector<Measurement> read_measurements(const std::string& filename) {
//std::ifstream is used to read the .csv file and create an input file stream named file
//where {filename} initializes the stream
std::ifstream file{filename};
//creates empty measurements vector where objects from Measurement will be stored
std::vector<Measurement> measurements;
//define variables needed
int year, month, day;
int hour, minute, second;
double temp;
char quality;
char dash;
char colon;
//Extract values from the file stream in a while loop
// example: file >> variable will extract data from file and store the data in variable
// In this case it expects to read: year-month-day hour:minute:second temp quality
while (file >> year >> dash >> month >> dash >> day
>> hour >> colon >> minute >> colon >> second >> temp >> quality) {
//create object measurement which is used to test if the measurement is valid
//using the .is_valid member function that returns true if it is valid
Measurement measurement{year, month, day, hour, minute, second, temp, quality};
if (measurement.is_valid()) {
//.push_back member function is to add to the vector with evey iteration
measurements.push_back(measurement);
} else {
std::cerr << "Warning: Skipping invalid measurement: ("
<< year << "-" << month << "-" << day << ","
<< hour << ":" << minute << ":" << second << ","
<< temp << "," << quality << ")"
<< std::endl;
//std::endl is a stream manipulator that inserts a new line
//(used in place of /n)
//and flushes the output stream
//(meaning it forces data in a memory buffer to move)
}
}
return measurements;
}
//save the measurements to a ROOT file
//having the output file pass into the function allows to reuse this script
//with different output files
void persist_measurements(const std::vector<Measurement>& measurements,
const std::string& output_filename) {
//create ROOT Tfile
//c_str() is a function that returns a pointer (similar to using: const char*)
TFile output_file{output_filename.c_str(), "RECREATE"};
//create TTree with the name "MNXB11" and create datasets with a branch for each variable
TTree tree{"MNXB11", "MNXB11"};
Double_t temp;
Int_t year;
Int_t month;
Int_t day;
Int_t hour;
Int_t minute;
Int_t second;
Int_t quality;
tree.Branch("temp", &temp);
tree.Branch("year", &year);
tree.Branch("month", &month);
tree.Branch("day", &day);
tree.Branch("hour", &hour);
tree.Branch("minute", &minute);
tree.Branch("second", &second);
tree.Branch("quality", &quality);
//auto& is used to make a reference whos data type is determined by the compiler
for (const auto& measurement : measurements) {
year = measurement.get_year();
month = measurement.get_month();
day = measurement.get_day();
hour = measurement.get_hour();
minute = measurement.get_minute();
second = measurement.get_second();
temp = measurement.get_temp();
quality = measurement.get_quality();
tree.Fill();
}
tree.Write();
}

Within the DataExtraction.cxx file, there are notes to explain each portion of the script. As an
overview, it starts by using the std::ifstream class to read files and create a vector using std::vector.
The while loop extracts the data from the data file in the expected format and order that is shown in
the parameters for the while loop. Using the .push_back member function it adds a Measurement
object to the empty vector. Persist_measurements is created to save the measurements to an
output file. The ROOT file is created using the output file and "RECREATE" with a tree named
"MNXB11". Each variable is then connected to its own branch in the ROOT tree by defining them
as ROOT variables and connecting them each to a branch using tree.Branch(). The for loop at the
end is used to get the stored values using the getter functions defined in Measurement.cxx. The
values are stored with tree.Fill() and the tree.Write writes the tree to the file.
